\begin{definition}[Normal rectangular injectivity hypotheses]%
    \label{def:peps_normalRectangleInjectivityHypotheses}
    \lean{TNLean.PEPS.NormalRectangleInjectivityHypotheses}
    \lean{TNLean.PEPS.NormalSquareLatticeRectangleInjectivityHypotheses}
    \leanok
    \uses{def:peps_regionInjectivityData,
        def:peps_squareLatticeCoordinateRegions}
    In the square-lattice normal theorem, every $2\times 3$ and
    $3\times 2$ rectangular region is assumed injective. The corresponding
    abstract hypotheses specify which finite regions have these two shapes and
    assert injectivity for all of them. In the coordinate square-lattice
    specialization, these two families are the contiguous coordinate
    $2\times 3$ and $3\times 2$ rectangles.
\end{definition}

\begin{lemma}[Coordinate rectangular injectivity gives abstract rectangular
        injectivity]%
    \label{lem:peps_squareLatticeRectangleInjectivity_toRectangular}
    \lean{TNLean.PEPS.NormalSquareLatticeRectangleInjectivityHypotheses.toRectangular}
    \leanok
    \uses{def:peps_normalRectangleInjectivityHypotheses}
    The coordinate square-lattice rectangular injectivity hypotheses determine
    abstract rectangular injectivity hypotheses, with the two abstract
    rectangular families taken to be the contiguous coordinate $2\times 3$
    and $3\times 2$ rectangles.
\end{lemma}
\begin{proof}\leanok
    This is the direct specialization of the abstract rectangular families to
    the two contiguous coordinate-rectangle predicates.
\end{proof}

\begin{lemma}[Coordinate rectangular injectivity gives injective rectangles]%
    \label{lem:peps_squareLatticeRectangleInjectivity_rectangles}
    \lean{TNLean.PEPS.NormalSquareLatticeRectangleInjectivityHypotheses.rect23_injective}
    \lean{TNLean.PEPS.NormalSquareLatticeRectangleInjectivityHypotheses.rect32_injective}
    \leanok
    \uses{def:peps_normalRectangleInjectivityHypotheses,
        lem:peps_squareLatticeCoordinateRegion_identities}
    Under the coordinate square-lattice rectangular injectivity hypotheses,
    every bounded contiguous $2\times3$ rectangle and every bounded
    contiguous $3\times2$ rectangle is injective.
\end{lemma}
\begin{proof}\leanok
    This is the defining rectangular injectivity hypothesis applied to the
    coordinate witnesses for the two rectangle shapes.
\end{proof}

\begin{lemma}[Rectangular injectivity gives injective $T$-edge blocks]%
    \label{lem:peps_normalSquareRegionT_edgeBlocks_injective}
    \lean{TNLean.PEPS.NormalSquareLatticeRectangleInjectivityHypotheses.tVerticalBlock_injective}
    \lean{TNLean.PEPS.NormalSquareLatticeRectangleInjectivityHypotheses.tHorizontalBlock_injective}
    \lean{TNLean.PEPS.NormalSquareLatticeRectangleInjectivityHypotheses.tHole_injective_of_union}
    \leanok
    \uses{def:peps_regionInjectivityData,
        def:peps_normalRectangleInjectivityHypotheses,
        lem:peps_normalSquareRegionTHole_rectangular_decomposition}
    For the displayed local $T$-region inside its $5\times6$ coordinate window,
    if the smaller edge-block bounds
    $x_0 + 5 \leq n$ and $y_0 + 5 \leq m$ hold in the finite $n\times m$
    square lattice, then the coordinate square-lattice rectangular injectivity
    hypotheses imply that the vertical $2\times3$ block and horizontal
    $3\times2$ block removed from the window are injective. If, in addition,
    region injectivity is closed under unions, then the union of these two
    removed edge blocks is injective.
\end{lemma}
% Maintainer note: this does not prove injectivity of the finite-lattice
% complementary block used as the third region in
% \cite[Theorem~3]{Molnar2018NormalPEPS}.
\begin{proof}\leanok
    Apply the rectangular injectivity hypotheses to the already identified
    contiguous rectangle shapes of the two edge blocks. The final assertion is
    one application of union closure to these two injective blocks.
\end{proof}

\begin{lemma}[Injectivity from a rectangular cover of the local $T$-region]%
    \label{lem:peps_normalSquareRegionT_injective_of_cover}
    \lean{TNLean.PEPS.NormalSquareLatticeRectangleInjectivityHypotheses.regionT_injective_of_cover}
    \lean{TNLean.PEPS.NormalSquareVerticalRegionTRectangleCover}
    \lean{TNLean.PEPS.NormalSquareLatticeRectangleInjectivityHypotheses.verticalT_inj_of_cover}
    \leanok
    \uses{def:peps_normalSquareRectangleCover,
        def:peps_normalRectangleInjectivityHypotheses,
        lem:peps_normalSquareRectangleCover_injective,
        lem:peps_regionInjectivityUnionClosure_finite}
    If the displayed local $T$-region has a nonempty finite cover by
    contiguous $2\times3$ and $3\times2$ rectangles, then rectangular
    injectivity and finite union closure imply that $T$ is injective. The
    rotated vertical local $T$-region has the same conditional injectivity
    criterion.
\end{lemma}
% Current source-comparison status: docs/paper-gaps/peps_normal_ft_section3_route.tex.
\begin{proof}\leanok
    Apply rectangular injectivity to every rectangle in the cover, then apply
    finite union closure to the nonempty family.
\end{proof}

\begin{lemma}[Injectivity from a rectangular cover of the edge complement]%
    \label{lem:peps_normalSquareEdgeComplementRegion_injective_of_cover}
    \lean{TNLean.PEPS.NormalSquareLatticeRectangleInjectivityHypotheses.edgeComplement_injective}
    \lean{TNLean.PEPS.NormalSquareLatticeRectangleInjectivityHypotheses.%
        verticalEdgeComplement_injective}
    \lean{TNLean.PEPS.NormalSquareLatticeRectangleInjectivityHypotheses.%
        edgeComplement_injective_of_T_cover}
    \leanok
    \uses{def:peps_normalSquareRectangleCover,
        def:peps_normalRectangleInjectivityHypotheses,
        lem:peps_normalSquareRectangleCover_injective,
        lem:peps_regionInjectivityUnionClosure_finite,
        lem:peps_normalSquareEdgeComplementRegion_coverOfT}
    If the finite-lattice complementary block $A_3$ has a nonempty finite
    cover by contiguous $2\times3$ and $3\times2$ rectangles, then
    rectangular injectivity and finite union closure imply that $A_3$ is
    injective. The same criterion applies to every vertical edge-complementary
    block: a nonempty finite cover of that block by contiguous $2\times3$ and
    $3\times2$ rectangles implies injectivity of the vertical
    edge-complementary block. In the normalized horizontal-edge $5\times7$
    frame, it is enough to give such a cover of the local $T$-region, since
    adjoining the two top-collar $3\times2$ rectangles gives a cover of $A_3$.
\end{lemma}
% Current source-comparison status: docs/paper-gaps/peps_normal_ft_section3_route.tex.
\begin{proof}\leanok
    Apply rectangular injectivity to every rectangle in the cover, then apply
    finite union closure to the nonempty family. For the normalized
    $5\times7$ edge complement, first extend the local $T$-cover by the
    two top-collar rectangles.
\end{proof}

\begin{lemma}[A cover of $T$ extends to a cover of the edge complement]%
    \label{lem:peps_normalSquareEdgeComplementRegion_coverOfT}
    \lean{TNLean.PEPS.normalSquareEdgeComplementRectangleCoverOfT}
    \lean{TNLean.PEPS.normalSquareVerticalEdgeComplementRectangleCoverOfVerticalT}
    \leanok
    \uses{def:peps_normalSquareRectangleCover,
        lem:peps_normalSquareEdgeComplementRegion_topCollar}
    In the normalized horizontal-edge $5\times7$ frame with lower-left corner
    $(0,0)$, any rectangular cover of the local $T$-region extends to a
    rectangular cover of the edge-complementary block $A_3$ by adjoining the
    two top-collar contiguous $3\times2$ rectangles. In the normalized
    vertical-edge $7\times5$ frame with lower-left corner $(0,0)$, any
    rectangular cover of the rotated local $T$-region extends to a rectangular
    cover of the vertical edge-complementary block by adjoining the two
    right-collar contiguous $2\times3$ rectangles.
\end{lemma}
\begin{proof}\leanok
    Index the new cover by the old cover indices together with two extra
    indices for the collar rectangles. For the horizontal block, the
    decomposition is
    \[
        A_3 =
        T \cup (\{0,1,2\}\times\{5,6\})
        \cup (\{2,3,4\}\times\{5,6\}).
    \]
    For the vertical block, the decomposition is
    \[
        A_3^{\mathrm v} =
        T^{\mathrm v} \cup (\{5,6\}\times\{0,1,2\})
        \cup (\{5,6\}\times\{2,3,4\}).
    \]
\end{proof}

\begin{lemma}[The normalized edge complements have collars]%
    \label{lem:peps_normalSquareEdgeComplementRegion_topCollar}
    \lean{TNLean.PEPS.normalSquareEdgeComplementRegion_eq_T_union_topCollar}
    \lean{TNLean.PEPS.normalSquareEdgeComplementRegion_eq_T_union_rightCollar}
    \lean{TNLean.PEPS.normalSquareVerticalRegionTHole}
    \lean{TNLean.PEPS.normalSquareVerticalRegionT}
    \lean{TNLean.PEPS.mem_normalSquareVerticalRegionTHole}
    \lean{TNLean.PEPS.mem_normalSquareVerticalRegionT}
    \lean{TNLean.PEPS.normalSquareVerticalEdgeComplement_eq_verticalT_union_rightCollar}
    \lean{TNLean.PEPS.NormalSquareLatticeRectangleInjectivityHypotheses.topCollar_injective}
    \lean{TNLean.PEPS.NormalSquareLatticeRectangleInjectivityHypotheses.rightCollar_injective}
    \lean{TNLean.PEPS.NormalSquareLatticeRectangleInjectivityHypotheses.verticalComp_injective}
    \lean{TNLean.PEPS.NormalSquareLatticeRectangleInjectivityHypotheses.verticalComp_inj_of_cover}
    \leanok
    \uses{lem:peps_normalSquareEdgeComplementRegion,
        lem:peps_normalSquareRegionT_window_complement,
        def:peps_regionInjectivityData,
        lem:peps_squareLatticeRectangleInjectivity_rectangles}
    In the normalized horizontal-edge $5\times7$ frame, the finite-lattice
    complementary block $A_3$ is the local-window $T$-region together with
    two top-collar contiguous $3\times2$ rectangles. Consequently, if the
    local $T$-region is injective, then rectangular injectivity and union
    closure imply that $A_3$ is injective. There is also a transposed
    $7\times5$ horizontal-frame collar identity with two contiguous
    $2\times3$ right-collar rectangles. For the rotated vertical
    edge-blocking regions, the corresponding complement is the rotated local
    $T$-region together with the same two right-collar rectangles; hence rotated
    $T$-injectivity, and hence a rectangular cover of rotated $T$, implies
    injectivity of the vertical complementary block.
    \begin{tenkzequation}
        % Formula verdict: in the normalized 7-by-5 row-column frame,
        % A3=((1-7,1-5)-(3-5,1-2)-(2-3,3-5)) is exactly the local T
        % region together with collars C1=(6-7,1-3) and C2=(6-7,3-5).
        % Ink-to-index verdict: every dot is an ordinary tensor and every
        % lattice edge is a virtual contraction; region fills encode only the
        % stated cell-set equality and do not alter the contraction graph.
        % Contraction/boundary verdict: each 7-by-5 panel has 58 internal
        % virtual contractions and 24 cut virtual bonds, with no exposed
        % physical indices, hence boundary type (V,P)=(24,0).
        % Source verdict: arXiv:1804.04964, paper_normal.tex:1430--1445 and
        % 1475--1499 motivate the complementary block.  The collar equality is
        % the exact local formal decomposition in NormalEdgeComplementCover.lean.
        \begin{tenkzlattice}[rows=7, cols=5, boundary legs]
            \tnregion[slot=selected, label=T, label at=north]
                {(1-6,1-5) - (3-5,1-2) - (2-3,3-5)}
            \tnregion[slot=collar, outline, label={C_1},
                label at=center]{(6-7,1-3)}
            \tnregion[slot=collar, outline, inset=1, label={C_2},
                label at=center]{(6-7,3-5)}
        \end{tenkzlattice}
        $\;=\;$
        \begin{tenkzlattice}[rows=7, cols=5, boundary legs]
            \tnregion[slot=complement, label={A_3}, label at=north]
                {(1-7,1-5) - (3-5,1-2) - (2-3,3-5)}
        \end{tenkzlattice}
    \end{tenkzequation}
\end{lemma}
% Current source-comparison status: docs/paper-gaps/peps_normal_ft_section3_route.tex.
\begin{proof}\leanok
    The two set identities follow by expanding the coordinate inequalities.
    The injectivity statements apply union closure first to the local
    $T$-region and the first collar rectangle, and then to the second collar
    rectangle.
\end{proof}

\begin{lemma}[A large rectangle is injective]%
    \label{lem:peps_normalSquareWideRectangle_injective}
    \lean{TNLean.PEPS.NormalSquareLatticeRectangleInjectivityHypotheses.wideRectangle_injective}
    \lean{TNLean.PEPS.NormalSquareLatticeRectangleInjectivityHypotheses.shortRectangle_injective}
    \leanok
    \uses{lem:peps_injectiveRegion_union,
        lem:peps_squareLatticeRectangleInjectivity_rectangles}
    A contiguous rectangle whose width is at least two and whose height is at
    least three is the union of the source $2\times3$ rectangles sliding over
    it, so rectangular injectivity and union closure prove it injective. The
    transposed statement holds for a contiguous rectangle of width at least
    three and height at least two.
\end{lemma}
% Current source-comparison status: docs/paper-gaps/peps_normal_ft_section3_route.tex.
\begin{proof}\leanok
    Each cell of the rectangle lies in one of the contained source windows, so
    the rectangle equals the finite union of those windows; union closure then
    gives injectivity.
\end{proof}

\begin{lemma}[Interior edge-complement injectivity]%
    \label{lem:peps_normalSquareInteriorEdgeComplement_injective}
    \lean{TNLean.PEPS.NormalSquareLatticeRectangleInjectivityHypotheses.interiorEdgeComplement_injective}
    \lean{TNLean.PEPS.NormalSquareLatticeRectangleInjectivityHypotheses.interiorVerticalEdgeComplement_injective}
    \leanok
    \uses{lem:peps_normalSquareEdgeComplementRegion,
        lem:peps_normalSquareWideRectangle_injective,
        lem:peps_squareLatticeRectangleInjectivity_rectangles}
    Let the two edge blocks be removed at an interior position of a sufficiently
    large lattice, with one row and column of room below and to the left and
    further room above and to the right. Then the complementary block $A_3$ is
    the union of four full-length bands around the removed blocks --- below,
    above, left, and right --- together with two filler rectangles inside the
    bounding box of the removed blocks. Each band is a large contiguous
    rectangle and each filler is a source rectangle, all avoiding the removed
    blocks, so rectangular injectivity and union closure prove $A_3$ injective
    with no rectangular-cover hypothesis. The transposed statement holds for the
    rotated vertical complementary block.

    \emph{Scope restriction (interior edge).} The removed blocks must sit in the
    interior. The fixed origin window has its corner forced and so has no
    rectangular cover; the source object is the full-lattice complement $A_3$,
    which has room only at an interior edge. The source comparison is recorded
    in \path{docs/paper-gaps/peps_normal_ft_section3_route.tex}.
\end{lemma}
% Current source-comparison status: docs/paper-gaps/peps_normal_ft_section3_route.tex.
\begin{proof}\leanok
    The six pieces cover the complement of the removed blocks, by expanding the
    coordinate inequalities. Each band is injective by the large-rectangle
    lemma and each filler by rectangular injectivity, so union closure gives
    injectivity of the union.
\end{proof}

\begin{lemma}[Cover-free interior every-edge blocking]%
    \label{lem:peps_normalSquareInteriorEdgeBlocking}
    \lean{TNLean.PEPS.horizontalSquareLatticeEdge_blockingDatum_interior}
    \lean{TNLean.PEPS.verticalSquareLatticeEdge_blockingDatum_interior}
    \lean{TNLean.PEPS.NormalSquareInteriorEdgeDatum}
    \lean{TNLean.PEPS.normalSquareEdgeBlockingHypotheses_of_interiorData}
    \leanok
    \uses{lem:peps_normalSquareInteriorEdgeComplement_injective,
        lem:peps_normalSquareHorizontalTranslatedEdge_blockingData,
        def:peps_normalEdgeBlockingHypotheses}
    For an interior horizontal or vertical edge with the corresponding interior
    margins, the translated edge-blocking datum is built with no rectangular
    cover: its complementary-block injectivity is the interior $A_3$
    injectivity. A family of these cover-free interior data over all edges
    assembles the normal edge-blocking hypotheses, with the injective chain,
    endpoint membership, pairwise disjointness, and coverage at every edge.

    \emph{Scope restriction (interior edge).} The construction supplies data
    only at edges with interior margins. Edges near the lattice boundary of the
    open rectangular model lack such margins; on a translation-invariant lattice
    every edge is interior. The source comparison is recorded in
    \path{docs/paper-gaps/peps_normal_ft_section3_route.tex}.
\end{lemma}
% Current source-comparison status: docs/paper-gaps/peps_normal_ft_section3_route.tex.
\begin{proof}\leanok
    Each interior edge is horizontal or vertical; the corresponding interior
    blocking datum discharges its complementary-block injectivity through the
    interior $A_3$ injectivity. Assembling these per-edge data over all edges
    gives the edge-blocking hypotheses.
\end{proof}

\begin{lemma}[Normalized edge-blocking data]%
    \label{lem:peps_normalSquareHorizontalEdge_blockingData}
    \lean{TNLean.PEPS.normalSquareHorizontalEdge}
    \lean{TNLean.PEPS.normalSquareHorizontalEdge_isHorizontal}
    \lean{TNLean.PEPS.normalSquareHorizontalEdgeRed}
    \lean{TNLean.PEPS.normalSquareHorizontalEdgeBlue}
    \lean{TNLean.PEPS.normalSquareHorizontalEdgeComplement}
    \lean{TNLean.PEPS.normalSquareHorizontalEdge_blockingDatum}
    \lean{TNLean.PEPS.normalSquareHorizontalEdge_blockingData}
    \lean{TNLean.PEPS.normalSquareHorizontalEdge_blockingData_of_TCover}
    \lean{TNLean.PEPS.normalSquareHorizontalEdge_injective_chain}
    \lean{TNLean.PEPS.normalSquareHorizontalEdge_blockingDatum_of_TCover}
    \lean{TNLean.PEPS.normalSquareHorizontalEdge_injective_chain_of_TCover}
    \lean{TNLean.PEPS.normalSquareVerticalEdge}
    \lean{TNLean.PEPS.normalSquareVerticalEdge_isVertical}
    \lean{TNLean.PEPS.normalSquareVerticalEdgeRed}
    \lean{TNLean.PEPS.normalSquareVerticalEdgeBlue}
    \lean{TNLean.PEPS.normalSquareVerticalEdgeComplement}
    \lean{TNLean.PEPS.normalSquareVerticalEdge_blockingDatum}
    \lean{TNLean.PEPS.normalSquareVerticalEdge_blockingData}
    \lean{TNLean.PEPS.normalSquareVerticalEdge_blockingData_of_verticalT}
    \lean{TNLean.PEPS.normalSquareVerticalEdge_blockingData_of_verticalTCover}
    \lean{TNLean.PEPS.normalSquareVerticalEdge_blockingDatum_of_verticalT}
    \lean{TNLean.PEPS.normalSquareVerticalEdge_blockingDatum_of_verticalTCover}
    \lean{TNLean.PEPS.normalSquareVerticalEdge_injective_chain}
    \lean{TNLean.PEPS.normalSquareVerticalEdge_injective_chain_of_verticalT}
    \lean{TNLean.PEPS.normalSquareVerticalEdge_injective_chain_of_verticalTCover}
    \leanok
    \uses{def:peps_squareLatticeCoordinateRegions,
        lem:peps_normalSquareEdgeComplementRegion,
        lem:peps_normalSquareRegionT_edgeBlocks_injective,
        lem:peps_normalSquareEdgeComplementRegion_topCollar}
    In the normalized $5\times7$ horizontal-edge frame, the red block is the
    vertical $2\times3$ block, the blue block is the horizontal $3\times2$
    block, and the complementary block is the finite-lattice edge complement
    $A_3$. The distinguished edge is a horizontal edge of the square-lattice
    graph. The left endpoint of the distinguished edge lies in the red block,
    the right endpoint lies in the blue block, the three regions are pairwise
    disjoint, and their union is the whole finite square lattice.
    If the local $T$-region is injective, rectangular injectivity and union
    closure give injectivity of all three blocks; in particular a rectangular
    cover of $T$ suffices. The normalized $7\times5$ vertical-edge frame
    has the analogous red and blue rectangular blocks, a distinguished vertical
    edge, and the same set-theoretic blocking data; at this stage its
    complementary-block injectivity is either kept as an explicit hypothesis or
    obtained from
    rotated $T$-injectivity, and therefore from a rectangular cover of
    rotated $T$, by the vertical collar lemma. In both cases the named
    blocking datum supplies the corresponding three injective regions for the
    edge-blocked chain.
\end{lemma}
% Maintainer note: horizontal and vertical complement injectivity are reduced
% to local and rotated T-covers, respectively. The translated every-edge
% construction remains a separate step.
\begin{proof}\leanok
    The endpoint memberships are coordinate checks. Rectangular injectivity
    gives the rectangular red and blue blocks. For the horizontal edge, the
    collar lemma gives the complement, and the rectangular cover criterion
    removes the direct $T$-injectivity input when such a cover is supplied.
    For the vertical edge, complement injectivity is an explicit input or is
    obtained from the rotated $T$-cover. The disjointness and cover
    statements are finite-set complement identities.
\end{proof}

\begin{lemma}[Translated horizontal edge-blocking data]%
    \label{lem:peps_normalSquareHorizontalTranslatedEdge_blockingData}
    \lean{TNLean.PEPS.normalSquareHorizontalTranslatedEdge}
    \lean{TNLean.PEPS.normalSquareHorizontalTranslatedEdgeRed}
    \lean{TNLean.PEPS.normalSquareHorizontalTranslatedEdgeBlue}
    \lean{TNLean.PEPS.normalSquareHorizontalTranslatedEdgeComplement}
    \lean{TNLean.PEPS.normalSquareHorizontalTranslatedEdge_isHorizontal}
    \lean{TNLean.PEPS.normalSquareHorizontalTranslatedEdge_eq_rightEdge}
    \lean{TNLean.PEPS.normalSquareHorizontalTranslatedEdge_blockingDatum}
    \lean{TNLean.PEPS.normalSquareHorizontalTranslatedEdge_injective_chain}
    \lean{TNLean.PEPS.normalSquareHorizontalTranslatedEdge_blockingDatum_of_complementCover}
    \lean{TNLean.PEPS.normalSquareHorizontalTranslatedEdge_injective_chain_of_complementCover}
    \lean{TNLean.PEPS.horizontalTranslatedEdge_zero}
    \lean{TNLean.PEPS.horizontalTranslatedRed_zero}
    \lean{TNLean.PEPS.horizontalTranslatedBlue_zero}
    \lean{TNLean.PEPS.horizontalTranslatedComp_zero}
    \leanok
    \uses{def:peps_squareLatticeCoordinateRegions,
        lem:peps_normalSquareEdgeComplementRegion,
        lem:peps_normalSquareEdgeComplementRegion_injective_of_cover}
    For any coordinate origin whose horizontal blocking window lies in the
    finite square lattice, there is a distinguished horizontal edge with the
    same relative position as in the normalized picture. The translated red
    region is the $2\times3$ block, the translated blue region is the
    $3\times2$ block, and the complementary region is the finite-lattice
    complement of their union. If this complementary region is injective,
    these three regions form the one-edge blocking datum; if it has a
    rectangular cover, rectangular injectivity and union closure supply that
    complementary injectivity. In both forms the datum supplies the three
    injective regions used in the edge-blocked chain. The distinguished edge
    is the corresponding coordinate right edge. At coordinate origin $0,0$,
    the translated edge and the three translated regions specialize to the
    normalized horizontal picture.
\end{lemma}
\begin{proof}\leanok
    The endpoint memberships and horizontal orientation are coordinate
    checks. Rectangular injectivity gives the red and blue blocks. The
    complementary block is supplied either as an input or from a rectangular
    cover, and the disjointness and covering identities follow from the
    finite-set complement formulas.
\end{proof}

\begin{lemma}[Translated vertical edge-blocking data]%
    \label{lem:peps_normalSquareVerticalTranslatedEdge_blockingData}
    \lean{TNLean.PEPS.normalSquareVerticalTranslatedEdge}
    \lean{TNLean.PEPS.normalSquareVerticalTranslatedEdgeRed}
    \lean{TNLean.PEPS.normalSquareVerticalTranslatedEdgeBlue}
    \lean{TNLean.PEPS.normalSquareVerticalTranslatedEdgeComplement}
    \lean{TNLean.PEPS.normalSquareVerticalTranslatedEdge_isVertical}
    \lean{TNLean.PEPS.normalSquareVerticalTranslatedEdge_eq_upEdge}
    \lean{TNLean.PEPS.normalSquareVerticalTranslatedEdge_blockingDatum}
    \lean{TNLean.PEPS.normalSquareVerticalTranslatedEdge_injective_chain}
    \lean{TNLean.PEPS.normalSquareVerticalTranslatedEdge_blockingDatum_of_complementCover}
    \lean{TNLean.PEPS.normalSquareVerticalTranslatedEdge_injective_chain_of_complementCover}
    \lean{TNLean.PEPS.verticalTranslatedEdge_zero}
    \lean{TNLean.PEPS.verticalTranslatedRed_zero}
    \lean{TNLean.PEPS.verticalTranslatedBlue_zero}
    \lean{TNLean.PEPS.verticalTranslatedComp_zero}
    \leanok
    \uses{def:peps_squareLatticeCoordinateRegions,
        lem:peps_normalSquareEdgeComplementRegion_injective_of_cover}
    For any coordinate origin whose vertical blocking window lies in the
    finite square lattice, there is a distinguished vertical edge with the same
    relative position as in the normalized vertical picture. The translated
    red region is the $3\times2$ block, the translated blue region is the
    $2\times3$ block, and the complementary region is the finite-lattice
    complement of their union. If this complementary region is injective,
    these three regions form the one-edge blocking datum; if it has a
    rectangular cover, rectangular injectivity and union closure supply that
    complementary injectivity. In both forms the datum supplies the three
    injective regions used in the edge-blocked chain. The distinguished edge
    is the corresponding coordinate upward edge. At coordinate origin $0,0$,
    the translated edge and the three translated regions specialize to the
    normalized vertical picture.
\end{lemma}
\begin{proof}\leanok
    The endpoint memberships and vertical orientation are coordinate checks.
    Rectangular injectivity gives the red and blue blocks. The complementary
    block is supplied either as an input or from a rectangular cover, and the
    disjointness and covering identities follow from the finite-set complement
    formulas.
\end{proof}
